\documentclass[10pt,a4paper]{amsart}
\usepackage{amsmath,amssymb,amsthm,graphicx}
% For commutative diagrams we use Xy-pic
\usepackage[all]{xy}

% We use 2cell for 2-commutative diagrams.
\xyoption{2cell}
% Theorem environments.
%
\theoremstyle{plain}
\newtheorem{theorem}[subsection]{Theorem}
\newtheorem{proposition}[subsection]{Proposition}
\newtheorem{lemma}[subsection]{Lemma}

\theoremstyle{definition}
\newtheorem{definition}[subsection]{Definition}
\newtheorem{example}[subsection]{Example}
\newtheorem{exercise}[subsection]{Exercise}
\newtheorem{situation}[subsection]{Situation}

\theoremstyle{remark}
\newtheorem{remark}[subsection]{Remark}
\newtheorem{remarks}[subsection]{Remarks}

\numberwithin{equation}{subsection}

% Macros
%
\def\lim{\mathop{\mathrm{lim}}\nolimits}
\def\colim{\mathop{\mathrm{colim}}\nolimits}
\def\Spec{\mathop{\mathrm{Spec}}}
\def\Hom{\mathop{\mathrm{Hom}}\nolimits}
\def\Ext{\mathop{\mathrm{Ext}}\nolimits}
\def\SheafHom{\mathop{\mathcal{H}\!\mathit{om}}\nolimits}
\def\SheafExt{\mathop{\mathcal{E}\!\mathit{xt}}\nolimits}
\def\Sch{\mathit{Sch}}
\def\Mor{\mathop{\mathrm{Mor}}\nolimits}
\def\Ob{\mathop{\mathrm{Ob}}\nolimits}
\def\Sh{\mathop{\mathit{Sh}}\nolimits}
\def\NL{\mathop{N\!L}\nolimits}
\def\CH{\mathop{\mathrm{CH}}\nolimits}
\def\proetale{{pro\text{-}\acute{e}tale}}
\def\etale{{\acute{e}tale}}
\def\QCoh{\mathit{QCoh}}
\def\Ker{\mathop{\mathrm{Ker}}}
\def\Im{\mathop{\mathrm{Im}}}
\def\Coker{\mathop{\mathrm{Coker}}}
\def\Coim{\mathop{\mathrm{Coim}}}

% Boxtimes
%
\DeclareMathSymbol{\boxtimes}{\mathbin}{AMSa}{"02}

%
% Macros for moduli stacks/spaces
%
\def\QCohstack{\mathcal{QC}\!\mathit{oh}}
\def\Cohstack{\mathcal{C}\!\mathit{oh}}
\def\Spacesstack{\mathcal{S}\!\mathit{paces}}
\def\Quotfunctor{\mathrm{Quot}}
\def\Hilbfunctor{\mathrm{Hilb}}
\def\Curvesstack{\mathcal{C}\!\mathit{urves}}
\def\Polarizedstack{\mathcal{P}\!\mathit{olarized}}
\def\Complexesstack{\mathcal{C}\!\mathit{omplexes}}
% \Pic is the operator that assigns to X its picard group, usage \Pic(X)
% \Picardstack_{X/B} denotes the Picard stack of X over B
% \Picardfunctor_{X/B} denotes the Picard functor of X over B
\def\Pic{\mathop{\mathrm{Pic}}\nolimits}
\def\Picardstack{\mathcal{P}\!\mathit{ic}}
\def\Picardfunctor{\mathrm{Pic}}
\def\Deformationcategory{\mathcal{D}\!\mathit{ef}}
\usepackage{cleveref}
\crefname{originallemma}{original lemma}{original lemmas}
\Crefname{originallemma}{Original lemma}{Original lemmas}
\makeatletter
\newlabel{derived-lemma-leray-acyclicity}{{015E}{source}{}{sourcepointer-1}{}}
\newlabel{derived-lemma-leray-acyclicity@cref}{{[originallemma][0][]015E}{[1][0][]source}{}{}{}}
\makeatother
\crefname{originalreference}{original reference}{original references}
\Crefname{originalreference}{Original reference}{Original references}
\makeatletter
\newlabel{lemma-V-C-all-n-etale-h}{{0EWF}{source}{}{sourcepointer-2}{}}
\newlabel{lemma-V-C-all-n-etale-h@cref}{{[originalreference][0][]0EWF}{[1][0][]source}{}{}{}}
\makeatother
\crefname{originalreference}{original reference}{original references}
\Crefname{originalreference}{Original reference}{Original references}
\makeatletter
\newlabel{lemma-cohomological-descent-etale}{{0DDI}{source}{}{sourcepointer-3}{}}
\newlabel{lemma-cohomological-descent-etale@cref}{{[originalreference][0][]0DDI}{[1][0][]source}{}{}{}}
\makeatother
\crefname{originalreference}{original reference}{original references}
\Crefname{originalreference}{Original reference}{Original references}
\makeatletter
\newlabel{lemma-describe-pullback-pi-h}{{0EW8}{source}{}{sourcepointer-4}{}}
\newlabel{lemma-describe-pullback-pi-h@cref}{{[originalreference][0][]0EW8}{[1][0][]source}{}{}{}}
\makeatother
\crefname{originallemma}{original lemma}{original lemmas}
\Crefname{originallemma}{Original lemma}{Original lemmas}
\makeatletter
\newlabel{sites-cohomology-lemma-relative-Leray}{{0734}{source}{}{sourcepointer-5}{}}
\newlabel{sites-cohomology-lemma-relative-Leray@cref}{{[originallemma][0][]0734}{[1][0][]source}{}{}{}}
\makeatother
\crefname{originallemma}{original lemma}{original lemmas}
\Crefname{originallemma}{Original lemma}{Original lemmas}
\makeatletter
\newlabel{sites-cohomology-lemma-torsion}{{0DD7}{source}{}{sourcepointer-6}{}}
\newlabel{sites-cohomology-lemma-torsion@cref}{{[originallemma][0][]0DD7}{[1][0][]source}{}{}{}}
\makeatother
\title{Cohomological descent for the comparison of the h-topology with the small etale topos}
\author{The Stacks Project Authors}
\date{stacks-project:etale-cohomology.tex:local-snapshot-20261011; GNU FDL 1.2 or later}
\begin{document}
\maketitle
\noindent\textit{Editorial note.} All delivered mathematical text is the author's own, copied from the pinned Stacks Project chapter file etale-cohomology.tex; no statement or proof step was rewritten, reordered or removed.\par\medskip
\noindent\textit{Editorial note.} The counted proof is the author's complete argument: the proof is written out in full: the reduction to a single sheaf, the rewriting of a\_\{X,*\}a\_X\textasciicircum{}\{-1\}F, and the degeneration of the relative Leray spectral sequence with the two vanishing inputs.\par\medskip
\noindent\textit{Editorial note.} The label lemma-cohomological-descent-etale-h names the counted statement, and the proof environment immediately following it in the source is delivered verbatim as its proof.\par\medskip
\noindent\textit{Editorial note.} The counted statement is carried by the author's own lemma environment in the chapter source.\par\medskip
\noindent\textit{Editorial note.} The statement and proof are copied from the pinned chapter file; the six cited inputs are kept as external pointers with their canonical Stacks tags.\par\medskip
\noindent\textit{Allowed standard prerequisites (editorial).}\begin{itemize}
\item Torsion complexes have torsion representatives. Cohomology on Sites, the lemma allowing a complex with torsion cohomology sheaves to be represented by a bounded below complex of torsion abelian sheaves. Supplies the representative F\textasciicircum{}bullet used in the reduction.
\item The description of the pullback in the h-topology and the V-C statement. Same chapter, the description of a\_X\textasciicircum{}\{-1\} through the small pullback, and the V-C statement giving the vanishing of the higher pushforwards of epsilon\_X. Supplies the identities a\_\{X,*\}a\_X\textasciicircum{}\{-1\}F = F and the vanishing of the higher pushforwards of the two factors.
\item Leray's acyclicity lemma and the relative Leray spectral sequence. Derived Categories, the acyclicity lemma, and Cohomology on Sites, the relative Leray spectral sequence for the composite a\_X = epsilon\_X o pi\_X. Supplies the reduction from a complex to a single sheaf and the vanishing of the higher direct images of a\_X.
\end{itemize}\medskip
\ifdefined\hypertarget\hypertarget{sourcepointer-1}{}\fi
\noindent\textit{Source reference.} Original lemma 015E: sources/stacks-chow/tags.txt, source lines 1505--1505. Prerequisite pointer invoked at proof line(s) 22789; it names a standard fact of a sibling chapter that the counted argument uses, and every step of the counted proof built on it is retained verbatim in the delivered body.\par\medskip
\ifdefined\hypertarget\hypertarget{sourcepointer-2}{}\fi
\noindent\textit{Source reference.} Original reference 0EWF: sources/stacks-etale-cohomology/pinned-original-source.tex, source lines 22718--22718. Prerequisite pointer invoked at proof line(s) 22801; it names a standard fact of the same chapter that the counted argument uses, and every step of the counted proof built on it is retained verbatim in the delivered body.\par\medskip
\ifdefined\hypertarget\hypertarget{sourcepointer-3}{}\fi
\noindent\textit{Source reference.} Original reference 0DDI: sources/stacks-etale-cohomology/pinned-original-source.tex, source lines 21308--21308. Prerequisite pointer invoked at proof line(s) 22804; it names a standard fact of the same chapter that the counted argument uses, and every step of the counted proof built on it is retained verbatim in the delivered body.\par\medskip
\ifdefined\hypertarget\hypertarget{sourcepointer-4}{}\fi
\noindent\textit{Source reference.} Original reference 0EW8: sources/stacks-etale-cohomology/pinned-original-source.tex, source lines 22554--22554. Prerequisite pointer invoked at proof line(s) 22795; it names a standard fact of the same chapter that the counted argument uses, and every step of the counted proof built on it is retained verbatim in the delivered body.\par\medskip
\ifdefined\hypertarget\hypertarget{sourcepointer-5}{}\fi
\noindent\textit{Source reference.} Original lemma 0734: sources/stacks-chow/tags.txt, source lines 9048--9048. Prerequisite pointer invoked at proof line(s) 22800; it names a standard fact of a sibling chapter that the counted argument uses, and every step of the counted proof built on it is retained verbatim in the delivered body.\par\medskip
\ifdefined\hypertarget\hypertarget{sourcepointer-6}{}\fi
\noindent\textit{Source reference.} Original lemma 0DD7: sources/stacks-chow/tags.txt, source lines 16757--16757. Prerequisite pointer invoked at proof line(s) 22784; it names a standard fact of a sibling chapter that the counted argument uses, and every step of the counted proof built on it is retained verbatim in the delivered body.\par\medskip
\begin{lemma}
\label{lemma-cohomological-descent-etale-h}
Let $X$ be a scheme. For $K \in D^+(X_\etale)$ with torsion cohomology
sheaves the map
$$
K \longrightarrow Ra_{X, *}a_X^{-1}K
$$
is an isomorphism with $a_X : \Sh((\Sch/X)_h) \to \Sh(X_\etale)$ as above.
\end{lemma}

\begin{proof}
We first reduce the statement to the case where
$K$ is given by a single abelian sheaf. Namely, represent $K$
by a bounded below complex $\mathcal{F}^\bullet$ of torsion
abelian sheaves. This is possible by Cohomology on Sites, Lemma
\ref{sites-cohomology-lemma-torsion}. By the case of a
sheaf we see that
$\mathcal{F}^n = a_{X, *} a_X^{-1} \mathcal{F}^n$
and that the sheaves $R^qa_{X, *}a_X^{-1}\mathcal{F}^n$
are zero for $q > 0$. By Leray's acyclicity lemma
(Derived Categories, Lemma \ref{derived-lemma-leray-acyclicity})
applied to $a_X^{-1}\mathcal{F}^\bullet$
and the functor $a_{X, *}$ we conclude. From now on assume $K = \mathcal{F}$
where $\mathcal{F}$ is a torsion abelian sheaf.

\medskip\noindent
By Lemma \ref{lemma-describe-pullback-pi-h} we have
$a_{X, *}a_X^{-1}\mathcal{F} = \mathcal{F}$. Thus it suffices to show that
$R^qa_{X, *}a_X^{-1}\mathcal{F} = 0$ for $q > 0$.
For this we can use $a_X = \epsilon_X \circ \pi_X$ and
the Leray spectral sequence
(Cohomology on Sites, Lemma \ref{sites-cohomology-lemma-relative-Leray}).
By Lemma \ref{lemma-V-C-all-n-etale-h}
we have $R^i\epsilon_{X, *}(a_X^{-1}\mathcal{F}) = 0$ for $i > 0$
and $\epsilon_{X, *}a_X^{-1}\mathcal{F} = \pi_X^{-1}\mathcal{F}$.
By Lemma \ref{lemma-cohomological-descent-etale} we have
$R^j\pi_{X, *}(\pi_X^{-1}\mathcal{F}) = 0$ for $j > 0$.
This concludes the proof.
\end{proof}

\end{document}
